\documentclass[11pt,a4paper]{article}
\usepackage[T1]{fontenc}
\usepackage{lmodern}
\usepackage[margin=28mm]{geometry}
\usepackage{amsmath,amssymb,amsthm,mathtools}
\usepackage{microtype}
\usepackage{enumitem}
\usepackage[hidelinks]{hyperref}

\newtheorem{theorem}{Theorem}[section]
\newtheorem{lemma}[theorem]{Lemma}
\newtheorem{proposition}[theorem]{Proposition}
\newtheorem{corollary}[theorem]{Corollary}
\theoremstyle{definition}
\newtheorem{definition}[theorem]{Definition}
\theoremstyle{remark}
\newtheorem{remark}[theorem]{Remark}

\newcommand{\UU}{\mathcal U}
\newcommand{\K}{\mathsf K}
\newcommand{\B}{\mathsf B}
\newcommand{\Cont}[2]{[#1\mathbin{\Rightarrow}#2]}
\newcommand{\step}[2]{(#1\mathbin{\Rightarrow}_{\!s}#2)}
\newcommand{\ev}[1]{\mathopen{[\![}#1\mathclose{]\!]}}
\newcommand{\lub}{\bigvee}
\newcommand{\N}{\mathbb N}
\newcommand{\TyRel}{\mathsf{Ty}}
\newcommand{\ElemRel}{\mathsf{Elem}}
\newcommand{\trunc}[1]{\left\lVert #1\right\rVert}
\newcommand{\appl}{\mathbin{\cdot}}
\DeclareMathOperator{\lfp}{lfp}

\title{Scott Domains and Continuous Function Spaces\\
  \large Small compact bases and an independent logical relation}
\author{}
\date{}

\begin{document}
\maketitle

\begin{abstract}
Scott-continuous maps between Scott domains with small compact bases form
a Scott domain: directed and bounded suprema are pointwise, and its compact
elements are exactly bounded finite joins of compact step functions. We
give a constructive proof that keeps the smallness condition explicit and
uses neither decidable order nor choice. We then use least fixed points
and a tower of finite-stage projectors on the Scott filter domain to construct
an independent logical relation. Its typed laws and fundamental lemma
cover dependent products, annotated abstraction and function
extensionality, and yield mere syntactic product inversion. An appendix
develops the
embedding--projection category, its equivalence with conditional upper
semilattices, and initial algebras for endofunctors preserving countable
chain colimits.
\end{abstract}

The function-space theorem is proved first because its compact basis and
pointwise completeness underlie the later Scott filter construction.
The fixed-point and projector-stage results then supply the semantic
projector used in the logical relation. The appendix records the
separate categorical theory of embedding--projection pairs.

\section{Scott domains and small bases}\label{sec:bases}

We reason constructively. Fix a universe $\UU$ of \emph{small} types and a
larger universe $\UU^+$. Carriers of domains may belong to $\UU^+$, whereas
orders are small proposition-valued relations. Unless otherwise stated,
all families whose suprema are considered have small index types.
Existential assertions express mere existence: they do not include a
chosen witness. These conventions are relevant to the proof; in a
classical set-theoretic reading one may simply fix a size bound throughout.

\begin{definition}
Let $D$ be a partially ordered set. An indexed family $(x_i)_{i\in I}$ is
\emph{directed} if $I$ is inhabited and, for every $i,j\in I$, there exists
$k\in I$ such that $x_i\leq x_k$ and $x_j\leq x_k$. It is \emph{bounded}
if there exists $u\in D$ such that $x_i\leq u$ for every $i\in I$.

The order is \emph{directed complete} if every small directed family has a
supremum, and \emph{bounded complete} if every small bounded family has a
supremum. An element $a\in D$ is \emph{compact} if, for every small directed
family,
\[
 a\leq\lub_{i\in I}x_i
 \quad\Longrightarrow\quad
 \exists i\in I\; a\leq x_i.
\]
\end{definition}

\begin{definition}\label{def:scott}
A \emph{Scott domain with a small compact basis} consists of a partially
ordered set $D$, a least element $\bot_D$, directed and bounded completeness,
and a small indexed family of compact elements
\[
 \beta_D:\B_D\longrightarrow D
\]
such that, for every $x\in D$, the family
\[
 \B_D(x)=\{a\in\B_D\mid\beta_D(a)\leq x\}
\]
is directed and has supremum $x$.
\end{definition}

For constructive and predicative domain theory with small bases, see
de Jong and Escard\'o~\cite{dejong-escardo}. The basis need not be
injectively indexed. We normally write $a$ for
$\beta_D(a)$ when an index $a\in\B_D$ is used as a domain element.
The least-upper-bound part of the basis condition can equivalently be
written
\begin{equation}\label{eq:density}
 \bigl(\forall a\in\B_D\;(a\leq x\Longrightarrow a\leq y)\bigr)
 \quad\Longrightarrow\quad x\leq y.
\end{equation}
Our completeness convention is for small indexed families, not arbitrary
families indexed by types in $\UU^+$.

\begin{lemma}[Finite common bounds]\label{lem:finite}
If finitely many compact elements $a_1,\ldots,a_n$ lie below the supremum
of a directed family $(x_i)_{i\in I}$, then there exists $i\in I$ with
$a_r\leq x_i$ for every $r$. In particular, if $a_1,\ldots,a_n\leq x$
in a Scott domain, there exists $a\in\B_D(x)$ above all the $a_r$.
The statement includes $n=0$.
\end{lemma}
\begin{proof}
For each of the finitely many compact elements, compactness gives a
family member above it. Directedness combines two such witnesses; finite
induction combines all of them. For the empty family use inhabitedness
of $I$. For the second assertion apply the first to the directed basis
family with supremum $x$. Only finitely many existential witnesses are
combined, so no choice principle is used.
\end{proof}

\begin{lemma}\label{lem:compactjoins}
A bounded finite join of compact elements is compact. Every compact
element of a Scott domain is equal to some basis element, in the sense
of mere existence.
\end{lemma}
\begin{proof}
Suppose $a=\lub_{r=1}^n a_r$ and $a\leq\lub_i x_i$, where $(x_i)$ is
directed. Lemma~\ref{lem:finite} gives one $x_i$ above every $a_r$, hence
above $a$. This also proves compactness of the empty join $\bot$.
If $x$ is compact, apply its compactness to
$x=\lub\B_D(x)$. Some $a\in\B_D(x)$ satisfies $x\leq a$, and
antisymmetry gives $x=a$.
\end{proof}

\section{Continuous function spaces}\label{sec:function-space}

The proof proceeds from pointwise completeness to compact steps,
then to their bounded finite joins and a small basis.

\subsection{Pointwise suprema}

\begin{definition}
A function $f:D\to E$ is \emph{Scott-continuous} if it is monotone and,
for every small directed family $(x_i)$,
\[
 f\left(\lub_i x_i\right)=\lub_i f(x_i).
\]
Write $\Cont D E$ for the set of these functions, ordered pointwise.
\end{definition}

\begin{lemma}[A small order on functions]\label{lem:smallorder}
For $f,g\in\Cont D E$,
\[
 f\leq g
 \quad\Longleftrightarrow\quad
 \forall a\in\B_D\; f(a)\leq g(a).
\]
In particular the order on $\Cont D E$ can be expressed by a small
proposition, even when the carrier of $D$ is large.
\end{lemma}
\begin{proof}
Only the reverse implication needs proof. By continuity and the basis
property, for every $x$,
\[
 f(x)=\lub_{a\in\B_D(x)} f(a)
 \leq\lub_{a\in\B_D(x)} g(a)=g(x).
\]
\end{proof}

\begin{proposition}[Pointwise suprema]\label{prop:pointwise}
The poset $\Cont D E$ has a least element and all small directed and
bounded suprema. Its least element is the constant function $\bot_E$,
and these suprema are computed pointwise.
\end{proposition}
\begin{proof}
Let $(f_i)_{i\in I}$ be a directed or bounded family of continuous maps.
For every $x$, the corresponding family $(f_i(x))$ is directed or bounded,
respectively, so define
\[
 h(x)=\lub_{i\in I}f_i(x).
\]
Monotonicity of $h$ follows from that of every $f_i$ and the universal
property of the supremum.

To prove continuity, let $(x_j)_{j\in J}$ be directed, with supremum $x$.
Monotonicity already gives $h(x_j)\leq h(x)$. Suppose that $z$ is any upper
bound of all the $h(x_j)$. For every $i,j$ we have
$f_i(x_j)\leq h(x_j)\leq z$. Continuity of $f_i$ gives $f_i(x)\leq z$,
and therefore $h(x)\leq z$. Thus $h(x)$ is the supremum of the family
$(h(x_j))$.

The pointwise universal property immediately makes $h$ the supremum of
the $f_i$ in $\Cont D E$. The constant bottom function is continuous
and lies below every function. Taking the empty bounded supremum gives
this same function.
\end{proof}

\begin{lemma}[Finite input witnesses]\label{lem:witness}
If $f\in\Cont D E$, $x\in D$, and $b$ is compact in $E$, then
\[
 b\leq f(x)
 \quad\Longrightarrow\quad
 \exists a\in\B_D(x)\; b\leq f(a).
\]
\end{lemma}
\begin{proof}
Apply compactness of $b$ to the directed supremum
$f(x)=\lub_{a\in\B_D(x)}f(a)$.
\end{proof}

\subsection{Compact step functions}

Let $a$ and $b$ be compact elements of $D$ and $E$, respectively.
For each $x\in D$, the proposition $a\leq x$ is a small index type.
The constant $b$ family indexed by this proposition is bounded by $b$.
We may therefore define
\begin{equation}\label{eq:step}
 \step a b(x)=\lub_{p:\,a\leq x} b.
\end{equation}
This definition does not decide whether $a\leq x$. If $a\leq x$, its
value is $b$; if $\neg(a\leq x)$, its value is $\bot_E$. These conditional
statements do not assert a disjunction between the two cases.
The subsingleton-supremum presentation of steps is also used by
de Jong and Escard\'o~\cite{dejong-escardo}.

\begin{proposition}\label{prop:step}
The function $\step a b$ is continuous and compact in $\Cont D E$.
For every continuous $f:D\to E$,
\begin{equation}\label{eq:steptest}
 \step a b\leq f \quad\Longleftrightarrow\quad b\leq f(a).
\end{equation}
\end{proposition}
\begin{proof}
Monotonicity follows by enlarging the indexing proposition. To check
continuity on a directed family $(x_i)$ with supremum $x$, let $z$ bound
all the values $\step a b(x_i)$. If $a\leq x$, compactness of $a$ gives
an $i$ with $a\leq x_i$, whence $b\leq\step a b(x_i)\leq z$.
The supremum property in~\eqref{eq:step} now gives $\step a b(x)\leq z$.
The opposite inequalities follow by monotonicity.

If $\step a b\leq f$, evaluation at $a$ gives $b\leq f(a)$.
Conversely, if $b\leq f(a)$, then $a\leq x$ implies $b\leq f(x)$;
the supremum property gives $\step a b(x)\leq f(x)$.

Finally, suppose $\step a b\leq\lub_i f_i$ for a directed family of
continuous maps. By~\eqref{eq:steptest} and pointwise computation,
$b\leq\lub_i f_i(a)$. Compactness of $b$ gives an $i$ with
$b\leq f_i(a)$, and then $\step a b\leq f_i$.
\end{proof}

\subsection{Finite approximations}

The function-space argument can be given directly using bounded
completeness. Fix a continuous map $f:D\to E$. For compact basis elements
$a_i\in\B_D$ and $b_i\in\B_E$ satisfying $b_i\leq f(a_i)$, put
\[
 h=\bigvee_{i=1}^{m}\step{a_i}{b_i}.
\]
Each step lies below $f$ by Proposition~\ref{prop:step}. Thus this finite
family has the common upper bound $f$, and its supremum exists in the
function space by pointwise bounded completeness. The empty join is the
constant bottom function. There is no need to decide which entries are
active at an input, or to construct a function by evaluating a graph.

\begin{lemma}\label{lem:boundedsteps}
Every bounded finite join of compact steps is compact. Moreover,
\[
 \bigvee_{i=1}^{m}\step{a_i}{b_i}\leq g
 \quad\Longleftrightarrow\quad
 \bigwedge_{i=1}^{m} b_i\leq g(a_i).
\]
\end{lemma}
\begin{proof}
The characterization follows from the universal property of a supremum
and the step-function characterization. For compactness, suppose the
finite join lies below a directed supremum. Compactness of each step
gives an index witnessing its inclusion in the directed family. Finitely
many such witnesses have a common upper index, which then bounds their
join. The empty case uses inhabitedness of the directed family.
\end{proof}

\begin{definition}
For $f:D\to E$ continuous, let $\mathcal A_f$ be the indexed family
\[
 \mathcal A_f=
 \left\{
   \bigvee_{i=1}^{m}\step{a_i}{b_i}
   \;\middle|\;
   a_i\in\B_D,\ b_i\in\B_E,\ b_i\leq f(a_i)
 \right\}.
\]
The indices are finite lists of basis pairs with the displayed
inequalities; in particular this is a small family.
\end{definition}

\begin{lemma}\label{lem:stepdensity}
The family $\mathcal A_f$ is directed and has supremum $f$.
\end{lemma}
\begin{proof}
It contains the empty join. Combining two lists gives their join,
still bounded by $f$, so it is directed.

Let $h$ be its supremum. Clearly $h\leq f$. For the converse, fix $x\in D$
and a compact basis element $b\leq f(x)$. Continuity of $f$ and
algebraicity of $D$ give a compact basis element $a\leq x$ with
$b\leq f(a)$, by Lemma~\ref{lem:witness}. The singleton step belongs to
$\mathcal A_f$, and
\[
 b\leq\step a b(x)\leq h(x).
\]
Algebraicity of $E$ gives $f(x)\leq h(x)$.
\end{proof}

\subsection{The function-space theorem}

\begin{theorem}\label{thm:exponential}
If $D$ and $E$ are Scott domains, then $[D\Rightarrow E]$ is a Scott
domain. Its compact elements are precisely the bounded finite joins of
steps between compact basis elements.
\end{theorem}
\begin{proof}
The pointwise constructions give bottom, directed completeness and
bounded completeness. Lemmas~\ref{lem:boundedsteps} and
\ref{lem:stepdensity} give algebraicity. If $f$ is compact, apply its
compactness to the directed family $\mathcal A_f$, whose supremum is
$f$. Some member lies above $f$ and also below $f$, hence equals $f$.
Conversely, every member of every $\mathcal A_f$ is compact.

To obtain a single small compact basis, index finite lists of basis
pairs satisfying the following small boundedness condition: every finite
subfamily whose inputs have a common basis upper bound has outputs with
a common basis upper bound. Necessity follows by evaluating a bound of
the corresponding steps at a common input bound. Conversely, for each
$x$ take the directed supremum of joins of output subfamilies whose
inputs lie below $x$. Unions make this family directed, and its resulting
function is the join of the steps by their universal property. Thus the
lists satisfying the condition present all compact functions, and the
families $\mathcal A_f$ establish density.
\end{proof}

\begin{remark}[Local approximation and global smallness]
The family $\mathcal A_f$ is small for each given $f$. The final part of
the proof supplies one small basis for the entire function space, as
required by Definition~\ref{def:scott}.
\end{remark}

\begin{remark}[Joint boundedness]
For general Scott domains one must use joint boundedness of finite
families. Pairwise compatibility alone is insufficient. Classically, the
proper subsets of $\{1,2,3\}$ form a finite Scott domain under inclusion:
the three singleton subsets are pairwise bounded but have no common
upper bound. Taking a one-point input domain gives the corresponding
counterexample to a pairwise test on steps.
\end{remark}

\section{Least fixed points}

The following result requires only a pointed dcpo, although it applies
in particular to every Scott domain and, by
Theorem~\ref{thm:exponential}, to its continuous function spaces.

\begin{theorem}[Kleene]\label{thm:kleene}
Let $S$ be a pointed dcpo and $F:S\to S$ a Scott-continuous function.
Define
\[
 p_0=\bot_S,\qquad p_{n+1}=F(p_n),\qquad
 p=\lub_{n\in\N}p_n.
\]
Then $F(p)=p$. Moreover, for every $q\in S$,
\[
 F(q)\leq q\quad\Longrightarrow\quad p\leq q.
\]
Thus $p=\lfp(F)$ is the least fixed point, and is also below every
pre-fixed point of $F$.
\end{theorem}
\begin{proof}
Bottom gives $p_0\leq p_1$, and monotonicity of $F$ gives inductively
$p_n\leq p_{n+1}$. Hence the sequence is directed. Continuity gives
\[
 F(p)=\lub_n F(p_n)=\lub_n p_{n+1}=\lub_n p_n=p,
\]
where the last equality follows since the shifted chain is cofinal.
If $F(q)\leq q$, induction gives $p_0\leq q$ and
$p_{n+1}=F(p_n)\leq F(q)\leq q$. Taking the supremum proves $p\leq q$.
\end{proof}

In particular, after establishing continuity of a functional
$\Phi:\Cont D E\to\Cont D E$, one can define its least solution by
iteration from the constant bottom function. This construction assumes
no idempotence, deflation or finitariness property of its iterates or of
its limit. Any such properties must be proved separately.

\begin{theorem}[Admissible fixed-point induction]\label{thm:admissible-induction}
Let $F:S\to S$ be Scott-continuous on a pointed dcpo. Suppose a
predicate $\mathcal I$ is closed under small directed suprema, contains
$\bot$, and is preserved by $F$. Then
\[
 \mathcal I(\lfp(F)).
\]
No downward closure is required of $\mathcal I$.
\end{theorem}
\begin{proof}
Induction gives $\mathcal I(F^n(\bot))$ for every $n$; admissibility
gives the conclusion at their supremum. This argument is proved once
as a general rule. Its applications supply only admissibility, the
bottom case, and preservation by the defining functional.
\end{proof}

\path{ScottAdmissible.seg} checks this rule for every presented Scott
domain. The existing continuous binary-map interface has the same
rule in \path{KleeneConstruction.seg}, using the directed-supremum
closure interface of \path{BinaryAdmissible.seg}. Both rules allow
invariant evidence in the larger universe; neither requires resizing
it to a small proposition. They are derived theorems, not additional
axioms.

\section{Finite-stage projectors and projector compatibility}

We now consider the Scott domain $D$ of consistent filters generated by
finite tokens. The tokens are built from $\bot$, $U$, products
$\Pi(a,f)$, arrows $a\Rightarrow b$, and compatible intersections.
Domain elements are ordered by inclusion of their token observations.
Thus $\bot$ is least, and the intersection of compatible tokens
represents their join of information.
There are two different kinds of projector below. The binary semantic
projector $p(a,x)$ enforces typing; the unary $q_n$ observes a finite
depth of the domain equation. Only the latter supplies the locality
needed to construct the logical relation.

The filter domain comes with a fold--unfold isomorphism
\[
 \mathsf{unfold}:D\xrightarrow{\;\cong\;}\mathsf H(D),
 \qquad \mathsf{fold}:\mathsf H(D)\xrightarrow{\;\cong\;}D.
\]
Here $\mathsf H(D)$ records a universe-head proposition, an optional
product pair in $D\times\Cont D D$, and a function component in
$\Cont D D$. The universe and product heads are disjoint, and either
positive head forces the function component to bottom. For a
Scott-continuous map $q:D\to D$ with $q(\bot)=\bot$, let $\mathsf H[q]$
retain the universe head,
send a product pair $(a,g)$ to
$(q(a),q\circ g\circ q)$, and send the function component $f$ to
$q\circ f\circ q$. This is the projector action of the domain equation.

\begin{proposition}[Finite-stage projector tower]\label{prop:projector-tower}
Define $q_0$ by
\[
 q_0(x)=\lub\{k\in\{\bot,\mathbf u\}\mid k\leq x\},
\]
where $\mathbf u\in D$ is the compact element represented by the
universe token $U$. Define the remaining maps by the domain equation:
\[
 q_{n+1}=\mathsf{fold}\circ\mathsf H[q_n]\circ\mathsf{unfold}.
\]
These are Scott-continuous deflations and satisfy
\[
 q_n^2=q_n,\qquad q_nq_m=q_{\min(n,m)},\qquad
 q_n\leq q_{n+1}\leq\mathrm{id},\qquad
 \bigvee_n q_n=\mathrm{id}.
\]
Put $D_n=\{x\in D\mid q_n(x)=x\}$, let $e_n:D_n\to D$ be inclusion, and
let $r_n:D\to D_n$ send $x$ to $q_n(x)$ with its fixed-point proof. Then
\[
 r_ne_n=\mathrm{id}_{D_n},\qquad e_nr_n=q_n\leq\mathrm{id}_D.
\]
Thus $q_0(x)$ has no positive product or arrow observation; its only
possible nonbottom compact observation is $\mathbf u$. Every compact
element is fixed by some $q_n$. On a positive product observation,
$q_{n+1}$ sends its domain component through $q_n$ and its family
function $g$ through $g\mapsto q_n\circ g\circ q_n$. It acts on an
ordinary function component by the same function projection and
preserves the universe observation. Consequently, components extracted
from an element of $D_{n+1}$ lie in, or factor through, $D_n$.
\end{proposition}
\begin{proof}
The indexing family defining $q_0(x)$ is inhabited by $\bot$,
directed, and bounded by $x$; in particular $q_0(\bot)=\bot$.
The head action preserves bottom, so every successive $q_n$ also
preserves bottom. For the successor, the displayed
fold--unfold equation applies $q_n$ to each domain component and
$g\mapsto q_n\circ g\circ q_n$ to each function component. It defines
a continuous map on filters independently of finite presentations.

Inductively, continuity, deflation and idempotence follow from the
corresponding properties of $q_n$; for the function component use
$q_n^2=q_n$ and monotonicity of application. The same componentwise
induction proves $q_n\leq q_{n+1}$ and the commuting equation. The
stage-zero image and the asserted successor support are immediate from
the defining clauses.

For exhaustivity, take a proper finite token presentation of a compact
element and proceed by structural induction. The tokens $\bot$ and $U$
are fixed by $q_0$. For an arrow token $a\Rightarrow b$, choose $n$
such that the compact elements represented by $a,b$ are fixed by
$q_n$; the function-component equation makes the arrow fixed by
$q_{n+1}$. For a product token $\Pi(a,f)$, the family token $f$ is
built from $\bot$, arrow tokens $a_i\Rightarrow b_i$, and finite
intersections. Let $F_f$ be the finite step function represented by
$f$. Choose one $n$ at which the compact elements represented by $a$
and all $a_i,b_i$ are fixed. Then $q_n\circ F_f\circ q_n=F_f$, so
the product-component equation fixes
$\Pi(a,f)$ at stage $n+1$. Finally, if two compatible compact elements
are fixed at stages $m$ and $n$, their join is fixed at
$\max(m,n)$: monotonicity puts both below the projection of their
join, while deflation gives the reverse inequality. Thus every compact
element is fixed by some $q_n$. Since these elements form a basis,
$\bigvee_n q_n(x)=x$ for every $x$.
\end{proof}

These $D_n$ are the fixed-image stages of the displayed fold--unfold
construction on the filter-domain solution. Their projections have not
been identified with those of the separate CUSL initial chain.

\subsection{Finitary projections through compact semilattices}

The general construction of $\mathrm{Fp}$ is particularly direct on the
CUSL side of the Scott--CUSL equivalence
(Theorem~\ref{thm:cusl-equivalence}). A CUSL is a partially ordered set
with bottom and joins of compatible pairs. Its ideal completion is a
Scott domain. The following construction uses its small presentation,
without requiring decidable compatibility or a resizing axiom.

\begin{theorem}[CUSL representation and the continuous $\mathrm{Fp}$ functor]
\label{thm:fp-cusl}
Let $C$ be a CUSL and $D=\operatorname{Idl}(C)$ its ideal completion.
Finitary projections on $D$ are represented by subsets $S\subseteq C$
containing bottom and closed under joins of pairs compatible in $C$:
\[
 p_S(x)=\lub\{c\in S\mid c\leq x\}.
\]
Such subsets are ordered by inclusion. They form an algebraic complete
lattice whose compact elements are the finitely generated closed
subsets. Their compact CUSL gives the CUSL presentation of
$\mathrm{Fp}(D)$.

For a CUSL embedding $e:C\hookrightarrow C'$, direct and inverse image
\[
 e_*(S)=e[S],\qquad e^*(T)=e^{-1}[T]
\]
form a continuous EP pair. This action satisfies the identity and
composition laws and preserves every existing colimit of a small
directed diagram of embeddings. Transport through the Scott--CUSL
equivalence gives the continuous functor $D\mapsto\mathrm{Fp}(D)$ on
Scott domains and EP pairs. For an EP pair $(e,r):D\hookrightarrow E$,
its embedding part is $P\mapsto ePr$.
\end{theorem}
\begin{proof}
For a closed subset $S$ and an ideal $x$, the elements of $S$ below $x$
are inhabited, directed and bounded. Their downward closure is the
ideal $p_S(x)$. This operation is continuous, deflationary and
idempotent. Its compact fixed points are exactly $S$, and they give
compact fixed-point density. Conversely, a finitary projection is
recovered from its compact fixed points. Thus the representation covers
all finitary projector operations, not only a specified subclass.
The subsets $S$ need not be downward closed: a compact element below a
fixed point need not itself be fixed.

Joins of closed subsets are obtained by closing their union under
compatible joins; directed joins are already unions. Closure of a
finite list of generators is compact. Every closed subset is the
directed union of these finitely generated subsets, so they form a
small compact basis. The checked compact-CUSL construction supplies
smallness of the resulting object.

An embedding preserves and reflects compatibility and preserves
bottom and compatible joins. Its direct images and inverse images
therefore remain closed. They satisfy
\[
 e^*e_*=\mathrm{id},\qquad e_*e^*(T)=T\cap e[C]\subseteq T.
\]
Both operations preserve directed unions. Their identity and
composition laws follow from those of images and inverse images.

For directed-colimit preservation, every element of the colimit CUSL
comes from a stage. This coverage is derived from the categorical
universal property: the union of the stage images inside the vertex
is itself a CUSL colimit, and uniqueness identifies its inclusion
with an isomorphism. Consequently every closed subset $T$ of the
vertex satisfies
\[
 T=\lub_i e_{i*}e_i^*(T).
\]
If $T$ is compact in the closed-family domain, it lies in one stage
image. Its preimage there is compact because EP embeddings reflect
compactness of elements in their images. Hence the mapped cocone also
has stage coverage, which gives its full universal property.
This proves preservation for arbitrary small directed index sets,
not just countable chains. The checked transport of colimits through
an equivalence gives the Scott-domain result.
\end{proof}

\begin{corollary}[The finitary-projector adjunction]\label{cor:fp-adjunction}
For every Scott domain $D$, inclusion
$i_D:\mathrm{Fp}(D)\hookrightarrow[D\to D]$ has a continuous right
adjoint $\kappa_D$, characterized by
\[
 i_D(P)\leq f\quad\Longleftrightarrow\quad P\leq\kappa_D(f).
\]
It satisfies $\kappa_D(i_D(P))=P$ and $i_D(\kappa_D(f))\leq f$;
thus $(i_D,\kappa_D)$ is an EP pair, as stated in
CGW~\cite[p.~162]{cgw}. In compact-family notation,
\[
 S_f=\{c\in K(D)\mid c\leq f(c)\},\qquad
 \kappa_D(f)(x)=\lub\{c\in S_f\mid c\leq x\}.
\]
The projection part of the functor in Theorem~\ref{thm:fp-cusl}, for
$(e,r):D\hookrightarrow E$, is
\[
 Q\longmapsto\kappa_D(rQe).
\]
\end{corollary}
\begin{proof}
The compact postfixed points $S_f$ contain bottom and are closed under
compatible joins. Their induced projector lies below $f$.
Compact fixed-point density shows that every finitary projector below
$f$ lies below this one, proving the adjunction and retraction.
A compact postfixed observation at a directed supremum of endomaps is
already observed at one index. Hence $\kappa_D$ preserves directed
suprema. The inclusion does too, since directed unions of closed
families give pointwise suprema of their projectors.
The general-domain adjunction is obtained from the ideal-domain one
by conjugating endomaps through the canonical ideal-completion
isomorphism. Finally, inverse image of compact fixed-point families
is exactly the family of compact postfixed points of $rQe$, proving
the displayed restriction formula.
\end{proof}

These results are checked independently of the filter domain and its
typing modules in \path{CUSLFpCheck.seg}. This session includes the
CUSL construction, the general directed-colimit theorem, its transport
to Scott domains, the adjunction and both action formulas. Source files
and individual theorem names are listed in \path{UNIT/README.txt}.

\subsection{The projector defined from the least fixed point}

Define an operator on the Scott domain of continuous binary maps,
\[
 \Phi:\Cont{D\times D}{D}\longrightarrow\Cont{D\times D}{D}.
\]
For $Q:\Cont{D\times D}{D}$, its value $\Phi(Q)$ is determined by the
following equations on the positive heads of the domain equation:
\[
\begin{aligned}
 \Phi(Q)(\mathbf u,\mathbf u)
   &=\mathbf u,\\
 \Phi(Q)(\mathbf u,\pi(a,f))
   &=\pi\!\left(Q(\mathbf u,a),
       z\mapsto Q\bigl(\mathbf u,f\appl Q(a,z)\bigr)\right),\\
 \Phi(Q)(\pi(a,f),x)
   &=\bigl(z\mapsto Q\bigl(f\appl Q(a,z),
                              x\appl Q(a,z)\bigr)\bigr).
\end{aligned}
\]
Here $f\appl y$ and $x\appl y$ denote semantic application, and the
last line denotes the corresponding function element of $D$.
All other combinations of basic heads give $\bot$. For arbitrary
elements, these equations are read through the constructive
fold--unfold decomposition of $D$: each positive head is
proposition-valued, and the three cases assemble into a continuous
binary map without deciding which head is present. This construction
is Scott-continuous in $Q$.
Define $p$ as the least fixed point of $\Phi$. By the
least-fixed-point theorem, it is computed by the Kleene sequence
\[
 P_0=\bot,\qquad P_{k+1}=\Phi(P_k),\qquad
 p=\lfp(\Phi)=\bigvee_{k\in\N}P_k.
\]
For any continuous binary operation $Q$, define its stage observation by
\[
 \rho_n(Q)(a,x)=r_n\bigl(Q(e_n a,e_n x)\bigr)
 \qquad(a,x\in D_n).
\]
At this point $p$ is a continuous binary operation satisfying
$\Phi(p)=p$. Its projector laws follow from the closure theorem and invariant below,
before comparison with inductive compact typing.

The dependent-product construction below follows Coquand, Gunter and
Winskel~\cite[Proposition~8 and \S7.2]{cgw}.

\begin{theorem}[Dependent products of subdomains]\label{thm:dependent-subdomain}
Let $A,D$ be Scott domains, and let $B:A\to\mathbf D^{\mathrm{EP}}$
be continuous, with coherent EP embeddings $B(a)\hookrightarrow D$.
Then the continuous sections form a Scott domain
\[
 \Gamma=\prod_{a\in A}B(a).
\]
If $A$ and $[D\to D]$ also embed in $D$ as EP pairs, then $\Gamma$
is a finitary subdomain of $D$.
For literal subdomains $A=\operatorname{Fix}(P)$ and
$B(a)=\operatorname{Fix}(Q_a)$, the intermediate EP pair is
\[
 E(s)(z)=s(Pz),\qquad R(g)(a)=Q_a(g(a)),
\]
and its projector on $[D\to D]$ is
\[
 (ER)(g)(z)=Q_{Pz}(g(Pz)).
\]
\end{theorem}
\begin{proof}
Represent a section by a continuous ambient-valued function $s:A\to D$
whose value at $a$ is fixed by the fiber projection $Q_a$.
The EP cone criterion converts continuity of $B$ into continuity of these
projections in the base argument. Conversely, a continuous family of
finitary projections gives a continuous EP family of image domains.
Both directions are checked using the explicit coherent inclusions and
the directed supremum of embedding--projection composites.

Directed and bounded joins of sections are pointwise. A compact input
$u\in A$ and a compact output $w\in D$ fixed by $Q_u$ determine a
compact section step. When it is active at $a$, monotonicity of the
family makes $w$ fixed by $Q_a$. Compatible finite joins of such steps
form a small compact basis. To cover a section observation, continuity
first gives a compact input witness, and finitariness of the corresponding
fiber projection gives a compact fixed output witness. This proves
algebraicity of $\Gamma$, without assuming an abstract section theorem.
The recovered fiber components also satisfy the monotonicity and
continuity equations for CGW sections.

The displayed $E$ and $R$ are continuous. Fiber fixedness and $P^2=P$
give $RE=\mathrm{id}$. For $g:D\to D$,
\[
 Q_{Pz}(g(Pz))\leq g(Pz)\leq g(z),
\]
so $ER\leq\mathrm{id}$. Composing this EP pair with $[D\to D]\hookrightarrow D$
embeds the Scott domain $\Gamma$ into $D$. Its induced projection is
finitary: embeddings preserve compactness, so the compact basis of
$\Gamma$ gives ambient compact fixed-point density.
The construction is checked for arbitrary presented Scott domains,
separately from the filter model and all typing modules.
\end{proof}

\begin{theorem}[Dependent sums and varying function spaces]\label{thm:sum-subdomain}
Let $A,C,D$ be Scott domains, let $a\mapsto Q_a$ be a continuous
family of finitary projections on $D$, and let $A\hookrightarrow C$
be an EP pair $(e,r)$. Then
\[
 \Sigma=\sum_{a\in A}\operatorname{Fix}(Q_a)
\]
is a Scott domain, ordered through its ambient components, with an EP
embedding into $C\times D$. Its induced projection is
\[
 (c,z)\longmapsto \bigl(e(rc),Q_{rc}(z)\bigr).
\]
If $R$ is a finitary projection on $D$, the varying spaces
\[
 V(a)=[\operatorname{Fix}(Q_a)\to\operatorname{Fix}(R)]
\]
have coherent embeddings into $[D\to D]$ and form a continuous EP
family. Their total space is therefore a finitary subdomain of
$C\times[D\to D]$, with projection
\[
 (c,g)\longmapsto
 \bigl(e(rc),\ z\mapsto R(g(Q_{rc}(z)))\bigr).
\]
\end{theorem}
\begin{proof}
Represent a point of the first sum by $(a,z)$ with $Q_a(z)=z$.
Directed and bounded joins are inherited from $A\times D$.
The compact basis consists of pairs $(u,w)$ where $u$ and $w$ are
ambient basis elements and $Q_u(w)=w$.
To prove density, consider compact demands $v\leq a$ and $b\leq z$.
Continuity of the family, applied to the basis approximation of $a$,
gives a compact base support $u_0\leq a$ with $b\leq Q_{u_0}(z)$.
Choose a compact $u\leq a$ above $v$ and $u_0$. Finitariness of $Q_u$
then gives a compact $w$ fixed by $Q_u$ with
$b\leq w\leq Q_u(z)\leq z$. Thus $(u,w)$ meets both demands.
This also gives directedness of the basis below each point.
This is the coherent-subdomain form of the support argument in
CGW~\cite[Proposition~7, pp.~132--134]{cgw}.

The embedding and projection are
$E(a,z)=(e(a),z)$ and $P(c,z)=(rc,Q_{rc}(z))$.
Joint continuity of $(a,z)\mapsto Q_a(z)$ proves continuity of $P$;
$PE=\mathrm{id}$ and $EP\leq\mathrm{id}$ follow from the fixed-point
condition and the given EP laws.

For the second assertion, the ordinary exponential EP action embeds
$V(a)$ into $[D\to D]$. Its induced finitary projection is
\[
 T_a(g)=R\circ g\circ Q_a.
\]
The exponential is order isomorphic to $\operatorname{Fix}(T_a)$.
The family $T_a$ is monotone and preserves directed suprema in $a$,
by continuity of $Q$, $g$ and $R$. Hence its fixed images form a
continuous EP family, and the first assertion applies.
All these constructions, including the compact bases and EP laws, are
checked for arbitrary presented Scott domains without filter or typing
modules. Compactness is used within these general domain constructions;
the application below uses their finitary projections.
\end{proof}

\begin{remark}[Evaluation over varying domains]\label{rem:varying-evaluation}
The same continuity mechanism appears in Coquand--Huber~\cite[p.~653
and Proposition~2, p.~658]{coquand-huber}. In their notation, consider
\[
 E=\sum_{a\in\mathrm{Type}}
 \bigl(\mathrm{El}(a)\times[\mathrm{El}(a)\to\mathrm{Type}]\bigr),
 \qquad \mathrm{ev}(a,v,f)=f(v).
\]
The evaluation map is continuous. Consequently, for compact
$w_0\leq f(v)$ there is a compact triple $(a_0,v_0,f_0)\leq(a,v,f)$
in this total space such that
\[
 w_0\leq f_0(v_0),\qquad
 v_0\in\mathrm{El}(a_0),\qquad
 f_0\in[\mathrm{El}(a_0)\to\mathrm{Type}].
\]
Here comparisons between different fibers use their coherent ambient
embeddings. The crucial point is that the compact argument and compact
function have the \emph{same} compact base code $a_0$.
Proposition~2 uses this simultaneous approximation to reduce an
observation of the codomain $f(v)$ to a finite product-type observation.

Theorem~\ref{thm:sum-subdomain} supplies the corresponding varying
function-space family and total-space support construction. Adding the
argument component gives the displayed space $E$. Joint continuity of
evaluation uses the same directed-diagonal argument as the continuous
projectors above: for a directed family of ambient representatives,
\[
 (\lub_i f_i)(\lub_i v_i)
   =\lub_{i,j}f_i(v_j)=\lub_k f_k(v_k).
\]
The last equality follows by choosing a common upper index for $i,j$.
Algebraicity of the total space and compactness of $w_0$ then give the
simultaneous witness. Thus the shared compact-base support principle
also connects these domain constructions to the logical relation.
The present Unit development checks the varying function spaces,
dependent-sum support construction and joint-continuity arguments used
above; the map $\mathrm{ev}$ on $E$ and its compact-witness corollary
are not yet packaged as separate checked theorems.
\end{remark}

\begin{theorem}[Normalized finitary families]\label{thm:normalized-finitary}
Suppose $Q:D\times D\to D$ is continuous, each $Q(a,-)$ is a finitary
projection, and
\[
 Q(Q(\mathbf u,a),x)=Q(a,x).
\]
Then $\Phi(Q)$ has the same properties. They hold for every Kleene
iterate $P_k$ and for $p$. In particular,
\[
 p(p(\mathbf u,a),x)=p(a,x).
\]
\end{theorem}
\begin{proof}
Apply admissible fixed-point induction, Theorem~\ref{thm:admissible-induction},
to the combined invariant
\[
 \mathcal I(Q)\quad\Longleftrightarrow\quad
 \bigl(\forall a,\;Q(a,-)\text{ is a finitary projection}\bigr)
 \;\wedge\;
 \bigl(\forall a,x,\;Q(Q(\mathbf u,a),x)=Q(a,x)\bigr).
\]
The induction hypothesis concerns the whole continuous family, so it
is available at every parameter occurring in a recursive call.

\emph{Preservation.}
Put $R=Q(a,-)$ and $S_v=Q(f(v),-)$. The function component is
\[
 T(g)(z)=S_{Rz}(g(Rz)).
\]
Apply Theorem~\ref{thm:dependent-subdomain} to the Scott domain
$\operatorname{Fix}(R)$ and the continuous EP family
$v\mapsto\operatorname{Fix}(S_v)$. The resulting subdomain projector,
after abstraction into $D$, is exactly the displayed function component.
The equality is checked in the filter instance of the general theorem.
This proves finitariness of the function component without normalization.

For the universe/product component put
$U_Q=\operatorname{Fix}(Q(\mathbf u,-))$ and
$D_a=\operatorname{Fix}(Q(a,-))$ for $a\in U_Q$.
Continuity of $Q$ gives a continuous family $a\mapsto D_a$.
By Theorem~\ref{thm:sum-subdomain}, the total space
\[
 \sum_{a\in U_Q}[D_a\to U_Q]
\]
is a finitary subdomain of $D\times[D\to D]$.
Its induced projection is
\[
 (a,f)\longmapsto
 \bigl(Q(\mathbf u,a),\ z\mapsto
 Q(\mathbf u,f(Q(Q(\mathbf u,a),z)))\bigr).
\]
Normalization identifies $Q(Q(\mathbf u,a),z)$ with $Q(a,z)$, so this
is precisely the universe/product payload operation of $\Phi$.
The equality is checked after abstraction in the filter model.
Its compact fixed-point density supplies the code and function
components simultaneously. Representing a compact function component
by a finite arrow code and applying the product constructor gives a
compact product code fixed by $\Phi(Q)(a,-)$.
The clause proof thus uses the general total-space theorem; it no longer
recursively assembles a common code from the entries of a finite graph.
The universe head itself is fixed; absent heads contribute bottom.
These cases give compact fixed-point coverage for $\Phi(Q)(a,-)$,
hence idempotence and finitariness.

For normalization, deflation gives one inequality. The reverse follows
by the observable-head decomposition. The product-type case uses
$Q(Q(\mathbf u,a),-)=Q(a,-)$ and idempotence of both $Q(a,-)$ and $Q(\mathbf u,-)$;
these identify its projected input and family in the function clause.
The universe and bottom cases follow directly from their heads.

\emph{Bottom.} The constant-bottom family satisfies both conditions.

\emph{Admissibility.} Let $(Q_i)$ be a small directed family satisfying
$\mathcal I$, and put $Q_\infty=\lub_iQ_i$. A compact fixed point of an
approximant remains fixed at the supremum,
so compact fixed-point density passes to the supremum. Deflation
passes pointwise, and compact fixed-point coverage gives idempotence.
For normalization, continuity
reduces the nontrivial inequality to mixed indices $n,m$. Choose
$i\geq n,m$ and use
\[
 Q_n(Q_m(\mathbf u,a),x)
 \leq Q_i(Q_i(\mathbf u,a),x)=Q_i(a,x)\leq Q_\infty(a,x).
\]
The reverse inequality follows from the diagonal terms.
Thus $\mathcal I$ is admissible. The induction rule now gives
$\mathcal I(p)$. This application contains no induction on the Kleene
iterates: \path{IndependentKleeneProjector.seg} supplies the three
obligations and invokes the rule. The independent projector checker
session checks that application and its preservation and admissibility
lemmas without the typed projector comparison.
\end{proof}

\begin{lemma}[One unfolding lowers the comparison stage]\label{lem:fusion}
For all continuous binary operations $Q,R$,
\[
 \rho_0(\Phi(Q))=\rho_0(\Phi(R)).
\]
For every $n$,
\[
 \rho_n(Q)=\rho_n(R)
 \quad\Longrightarrow\quad
 \rho_{n+1}(\Phi(Q))=\rho_{n+1}(\Phi(R)).
\]
\end{lemma}
\begin{proof}
At stage zero only the universe head can be observed, besides bottom.
The universe head of $\Phi(Q)(A,X)$ is present exactly when both $A$
and $X$ have a universe head; this does not depend on $Q$.

For the successor case, use the structural equations of
Proposition~\ref{prop:projector-tower}. On a product $q_{n+1}$ projects the domain with $q_n$ and the
family function by $g\mapsto q_n\circ g\circ q_n$. On the function
component it uses the same function projection, and it preserves the
universe head.

If $A,X$ belong to stage $n+1$, the domain component of either product
head belongs to stage $n$. Each family function, and the application
function of $X$, takes values in stage $n$ and factors through $q_n$.
Consequently, in the nested calls occurring in $\Phi(Q)$, a function
applied to $Q(e_n a,e_n z)$ depends only on
\[
 r_n Q(e_n a,e_n z)=\rho_n(Q)(a,z).
\]
Its output is itself an embedded stage-$n$ value. The outer recursive
call, after projection by $r_n$, is therefore also determined by
$\rho_n(Q)$.

For example, the restricted function clause is determined by
\[
 z\longmapsto
 \rho_n(Q)\bigl(B(\rho_n(Q)(a,z)),\,F(\rho_n(Q)(a,z))\bigr),
\]
where $a$ is the product's stage-$n$ domain component, $B$ its family,
and $F$ the stage-$n$ application function of $X$. The product clause
has the same calculation with the constant universe family. Replacing
$\rho_n(Q)$ by $\rho_n(R)$ gives equal components and hence equal
stage-$(n+1)$ results. These arguments use positive head observations;
no decidability of a head is required.
\end{proof}

\begin{corollary}[Finite stabilisation]\label{cor:stabilisation}
For $k,j\geq n+1$,
\[
 \rho_n(P_k)=\rho_n(P_j).
\]
\end{corollary}
\begin{proof}
Induct on $n$. The base is the first assertion of
Lemma~\ref{lem:fusion}. At the successor, write $k=k'+1$ and $j=j'+1$,
apply the induction hypothesis to $P_{k'},P_{j'}$, and then apply the
second assertion of the lemma.
\end{proof}

Define the stage operation directly from the least fixed point:
\[
 p_n(a,x)=r_n\bigl(p(e_n a,e_n x)\bigr),
 \qquad a,x\in D_n.
\]

\begin{theorem}[Finite computation of the stage operation]\label{thm:operational}
For every $a,x\in D_n$,
\[
 p_n(a,x)=r_n\bigl(P_{n+1}(e_n a,e_n x)\bigr).
\]
\end{theorem}
\begin{proof}
Continuity of $r_n$ and pointwise computation of the Kleene supremum give
\[
 \rho_n(p)=\bigvee_k\rho_n(P_k).
\]
The increasing family on the right is constant from index $n+1$ onwards,
by Corollary~\ref{cor:stabilisation}, so its supremum is
$\rho_n(P_{n+1})$, which is the displayed right-hand side.
\end{proof}

Thus $p_n$ can be computed using only $n+1$ unfoldings of $\Phi$.
This theorem does not prove that $p_n(a,-)$ is a finitary projection:
finite stabilisation alone gives neither deflation nor idempotence.
The checked stage finitariness proof used later identifies $p_n$ with
the projector generated by bounded inductive compact typing.
Theorem~\ref{thm:normalized-finitary} supplies independent global
finitariness, but the stage conclusion requires an additional argument.
Arbitrary EP conjugation $rPe$ need not be idempotent: on the finite
chain $0<1<2<3$, take $P=(0,0,2,2)$, embed $\{0,1,3\}$, and project
$2$ to $1$. Then $rPe$ sends $3$ to $1$ and $1$ to $0$.
The finite-order laws and this failure are checked in
\path{EPRestrictionObstruction.seg}. Compatibility with the particular
rank projections is therefore needed to remove the bounded-typing
stage proof.

The continuous functor $\mathrm{Fp}$ and its adjunction are now checked
in general (Theorem~\ref{thm:fp-cusl} and
Corollary~\ref{cor:fp-adjunction}). Its projection part sends a projector
$P$ to $\kappa_{D_n}(r_nPe_n)$. This produces a finitary projector,
but the direct stage formula above uses the raw restriction $r_nPe_n$.
To obtain an independent proof of that formula's finitariness, one must
show that these particular restrictions are unchanged by $\kappa_{D_n}$.
The existing bounded-typing proof establishes stage finitariness;
the general functor theorem does not replace the required rank
compatibility argument.


The logical-relation construction and fundamental lemma below use only
the projector-tower properties of Proposition~\ref{prop:projector-tower}.

\section{An independent logical relation}

\subsection{Admissible predicates and the relation clauses}

We now use the Scott filter domain to relate denotations to syntax.
Write $\mathbf u\in D$ for the universe element,
$\pi(a,\varphi)$ for the product constructor, and $f\appl x$ for
application. A family $\varphi$ is a function component of $D$, representing
a continuous map $D\to D$. The recursive domain equation supplies
universe observations and product observations. A positive product
observation
\[
 \pi(\bot,\bot)\leq x
\]
allows extraction of a domain component $x_{\rm d}$ and a family
component $x_{\rm f}$. These components are only used under that positive
hypothesis. No decision whether an arbitrary element has a head is needed.
The product $\pi(\bot,\bot)$ is a positive compact observation, distinct
from $\bot$.

Fix a syntactic context $\Delta$. Terms have variables, a universe $U$,
dependent products $\Pi(A,B)$, annotated abstractions $\lambda A.M$, and
applications. The body of a product or abstraction has one additional
variable; $B N$ denotes substitution of $N$ for it. More generally,
$A\sigma$ is substitution by the simultaneous substitution $\sigma$. Write
$T\longrightarrow^* V$ for head reduction and
$\Delta\vdash M=N:A$ for typed definitional equality. When such a
judgement occurs inside $\trunc{\;\cdot\;}$, it denotes the type of its
derivations, and the brackets denote propositional truncation.

\subsubsection{Lower admissible predicates}

\begin{definition}
A small proposition-valued predicate $Q$ on $D$ is \emph{lower admissible}
if it is downward closed and closed under small directed suprema.
It is \emph{strict} if $Q(\bot)$ holds. Strictness is a separate condition:
directed families are inhabited, so admissibility alone does not imply it.
\end{definition}

\begin{lemma}[Compact reconstruction]\label{lem:reconstruction}
For every lower admissible predicate $Q$,
\[
 Q(x)\quad\Longleftrightarrow\quad
 \forall k\in\B_D\;(k\leq x\Longrightarrow Q(k)).
\]
Conversely, for a small predicate $R$ on the compact basis,
\[
 \overline R(x)=\forall k\in\B_D\;(k\leq x\Longrightarrow R(k))
\]
is lower admissible. If $R$ respects the basis order, then
$\overline R(k)\Longleftrightarrow R(k)$ for every basis element $k$.
\end{lemma}
\begin{proof}
Downward closure proves one direction of reconstruction; the directed
basis of $x$ and admissibility prove the other. For admissibility of
$\overline R$, a compact observation below a directed supremum lies below
one member. Its mere witness may be eliminated into the proposition
$R(k)$. Recovery on a compact uses downward closure and the observation
of that compact itself.
\end{proof}

For lower admissible predicates $P,Q$, define preservation by
\[
 \operatorname{Arr}(P,Q)(f)
 \quad\Longleftrightarrow\quad
 \forall x\in D\;(P(x)\Longrightarrow Q(f\appl x)).
\]
Although this formula quantifies over a potentially large domain, it is
equivalent to the small formula
\begin{equation}\label{eq:compact-arrow}
 \forall k\in\B_D\;(P(k)\Longrightarrow Q(f\appl k)).
\end{equation}
Indeed, if $P(x)$, then $P(k)$ for every compact $k\leq x$; continuity
of application and admissibility of $Q$ give the conclusion at $x$.
For each compact input, application is continuous in $f$. Thus
\eqref{eq:compact-arrow} is an intersection of lower admissible predicates
of $f$, and is itself lower admissible. It is contravariant in $P$ and
covariant in $Q$.

\subsubsection{Type and term predicates}

We seek predicates
\[
 \TyRel(x;T,T'),\qquad \ElemRel(x;T,M,M')
\]
for all syntactic indices in $\Delta$. Each is to be small, strict and
lower admissible in $x$. There is no separate semantic classifier
argument in $\ElemRel$: semantic typing will enter its laws and the
fundamental lemma.

For candidate term predicates $P$ and $Q$, define the function condition
$\mathcal F(P,Q)(f;T,M,M')$ as follows. For every reduction
$T\longrightarrow^*\Pi(A,B)$, every derivation
$\Delta\vdash N=N':A$, and every $x$ with $P(x;A,N,N')$, require
\begin{equation}\label{eq:function-outputs}
\begin{aligned}
 Q(f\appl x;B N,MN,MN'),\\
 Q(f\appl x;B N,M'N,M'N'),\\
 Q(f\appl x;B N,MN,M'N').
\end{aligned}
\end{equation}
All three comparisons use the same left result type $B N$.
For candidate term predicates $P$ and type predicates $Q$, the family
condition $\mathcal G(P,Q)(\varphi;A,B,B')$ similarly requires
\begin{equation}\label{eq:family-outputs}
\begin{aligned}
 Q(\varphi\appl x;B N,B N'),\\
 Q(\varphi\appl x;B' N,B' N'),\\
 Q(\varphi\appl x;B N,B' N')
\end{aligned}
\end{equation}
whenever $\Delta\vdash N=N':A$ and $P(x;A,N,N')$.
The compact-input formulation makes both conditions small and lower
admissible in their semantic argument.

The type condition $\mathcal T(P,Q)(x;T,T')$ has two parts.
First, a universe observation $\mathbf u\leq x$ implies
\[
 \trunc{T\longrightarrow^*U\ \mathbin{\times}\ T'\longrightarrow^*U}.
\]
Second, a positive product observation implies the mere existence of
$A,A',B,B'$ with
\begin{equation}\label{eq:product-clause}
\begin{gathered}
 T\longrightarrow^*\Pi(A,B),\qquad
 T'\longrightarrow^*\Pi(A',B'),\\
 \Delta\vdash A=A':U,\qquad
 \Delta,A\vdash B=B':U,\\
 Q(x_{\rm d};A,A'),\qquad
 \mathcal G(P,Q)(x_{\rm f};A,B,B').
\end{gathered}
\end{equation}
The whole tuple in~\eqref{eq:product-clause} is truncated. In particular,
its component equalities are retained even at the smallest positive
product observation.

\begin{lemma}[Admissibility of the type condition]\label{lem:type-admissible}
The condition $\mathcal T(P,Q)$ is lower admissible whenever the
candidate predicates are lower admissible. It is contravariant in $P$
and covariant in $Q$. It is strict when the positive predicates are strict.
\end{lemma}
\begin{proof}
Head determinism says that any two evidenced product reducts of the same
pair $T,T'$ have identical syntactic components. Consequently the product
condition can equivalently be expressed by a mere existence of evidenced
heads, together with the component conditions for every such evidenced
head. This separates head existence from universal component tests.

Each head-existence requirement is an observation implication
$k\leq x\Longrightarrow E$, for a fixed compact head $k$ and a small
proposition $E$. Such an implication is lower admissible by compactness.
The domain component can be tested by observations $\pi(k,\bot)\leq x$;
the family component by observations $\pi(\bot,k)\leq x$. Compact
reconstruction identifies these tests with the required predicates on
the extracted components. Intersections give admissibility of the whole
condition. Variance follows from that of argument preservation. Bottom
has neither a universe head nor a positive product head, and the
remaining strict output conditions are inherited from $Q$.
\end{proof}

The desired invariant equations are
\begin{equation}\label{eq:relation-equations}
\begin{aligned}
 \TyRel(x;T,T')
  &\Longleftrightarrow\mathcal T(\ElemRel,\TyRel)(x;T,T'),\\
 \ElemRel(x;T,M,M')
  &\Longleftrightarrow\TyRel(x;M,M')\ \wedge\
    \mathcal F(\ElemRel,\ElemRel)(x;T,M,M').
\end{aligned}
\end{equation}
These equations are mixed in variance. In particular their simultaneous
diagonal action is not a monotone operator to which Kleene iteration
can simply be applied.

\subsection{Construction by compatible finite depths}

Use the projectors $q_d$ constructed above. They satisfy
$q_d\leq q_{d+1}\leq\operatorname{id}$,
$q_dq_e=q_{\min(d,e)}$, and $\lub_dq_d=\operatorname{id}$.
Every compact basis element is fixed by some $q_d$.
Two pairs of predicates \emph{agree at depth $d$} if their type and term
components are equivalent on every element fixed by $q_d$, for all
syntactic indices.

For a pair $R=(R^{\rm ty},R^{\rm tm})$, form the new type predicates
before the new term predicates:
\begin{equation}\label{eq:advance}
\begin{aligned}
 (R^+)^{\rm ty}(x;T,T')
  &=\mathcal T(R^{\rm tm},R^{\rm ty})(x;T,T'),\\
 (R^+)^{\rm tm}(x;T,M,M')
  &=(R^+)^{\rm ty}(x;M,M')\ \wedge\
    \mathcal F(R^{\rm tm},R^{\rm tm})(x;T,M,M').
\end{aligned}
\end{equation}
Reading the newly formed type component in the second line is essential.
The original mutual action reads an old type component at the same depth
and does not directly have the following one-step locality property.

\begin{lemma}[Locality]\label{lem:relation-locality}
If $R$ and $S$ agree at depth $d$, then $R^+$ and $S^+$ agree at
depth $d+1$.
\end{lemma}
\begin{proof}
For an element fixed by $q_{d+1}$, every extracted product domain is
fixed by $q_d$. Its family function, and its application function, have
depth-$d$ input and output support. Write $h$ for either such function.
Then
\[
 h(x)=h(q_dx),\qquad q_d(h(x))=h(x).
\]
Suppose the arrow condition for $R$ is known and an argument is related
by the negative predicate of $S$. Lower its evidence to $q_dx$.
Agreement at depth $d$ converts this evidence to the negative predicate
of $R$. Apply the given arrow condition at $q_dx$, use the first equation
to identify the output with $h(x)$, and use the second equation to
convert the positive output predicate from $R$ to $S$.

This proves both directions for functions and families; the domain
component is handled directly by agreement at depth $d$. Head evidence
is unchanged. Thus the new type predicates agree at depth $d+1$.
Using that agreement in the new term clause finishes the proof.
\end{proof}

Define the initial pair by the universe-observation tests
\[
\begin{aligned}
 R_0^{\rm ty}(x;T,T')&=
  \bigl(\mathbf u\leq x\Longrightarrow
    \trunc{T\longrightarrow^*U\mathbin{\times}T'\longrightarrow^*U}\bigr),\\
 R_0^{\rm tm}(x;T,M,M')&=R_0^{\rm ty}(x;M,M'),
 \qquad R_{d+1}=R_d^+.
\end{aligned}
\]
At depth zero, every element lies below $\mathbf u$, has no product head,
and its applications are bottom. Hence $R_0$ agrees at depth zero with
$R^+$ for every strict $R$. This statement does not require deciding
whether a depth-zero element is bottom or the universe.

\begin{lemma}[Compatibility and recovery]\label{lem:relation-recovery}
Each $R_d$ is strict and lower admissible, and
\[
 d\leq e,\quad q_dx=x
 \quad\Longrightarrow\quad R_d(x)\Longleftrightarrow R_e(x).
\]
Here and below such formulas apply separately to each component and
all of its syntactic indices. Define
\begin{equation}\label{eq:relation-limit}
 R_\infty(x)=\forall d\in\N\;R_d(q_dx).
\end{equation}
Then $R_\infty$ is strict and lower admissible, and
\[
 q_dx=x\quad\Longrightarrow\quad
 R_\infty(x)\Longleftrightarrow R_d(x).
\]
\end{lemma}
\begin{proof}
The depth-zero observation and Lemma~\ref{lem:relation-locality} prove
compatibility by induction. Strictness and admissibility are preserved
at every step. The limit is a small countable intersection of continuous
pullbacks of lower admissible predicates, so has the same properties.

For recovery, the forward implication uses the $d$th conjunct. Conversely
suppose $q_dx=x$ and $R_d(x)$. If $d\leq e$, compatibility gives $R_e(x)$,
which lowers to $R_e(q_ex)$. If $e\leq d$, first lower to $R_d(q_ex)$
and then use compatibility at the depth-$e$ element $q_ex$.
Constructive comparison of natural numbers gives every conjunct.
\end{proof}

\begin{theorem}[Independent invariant]\label{thm:independent-invariant}
There is a unique strict lower admissible pair $(\TyRel,\ElemRel)$ satisfying
\eqref{eq:relation-equations}. It is given by~\eqref{eq:relation-limit}.
\end{theorem}
\begin{proof}
Recovery and locality show that $R_\infty$ and $R_\infty^+$ agree at every
finite depth. Every compact is captured at a finite depth, so they agree
on all compact observations. Lemma~\ref{lem:reconstruction} gives agreement
on all domain elements. The type equation then identifies the fresh type
component in~\eqref{eq:advance} with $R_\infty^{\rm ty}$, yielding both
original mutual equations.

For uniqueness, a strict solution of the original equations also solves
the type-before-term equations. The depth-zero observation and locality
show inductively that it agrees with $R_d$ at depth $d$. Recovery and
compact reconstruction identify it with $R_\infty$.
\end{proof}

This is a finite-depth construction with an admissible limit, rather than
iteration of a monotone diagonal relation operator. It constructs the
relation without an earlier logical relation or adequacy theorem.
All existential eliminations have proposition-valued targets. In
particular, it never replaces $\forall i\,\trunc{W_i}$ by
$\trunc{\forall i\,W_i}$.

\subsection{Typed laws and restriction of arguments}

Write $v:c$ for the inductive compact typing judgement. For $a\in D$,
let $T(a)$ be the family generated by compact values $v$ for which
$v:c$ and $c\leq a$, closed under compatible finite joins and compact
equivalence. The associated map is
\[
 p_{\mathrm{ty}}(a,x)=\bigvee\{v\in T(a)\mid v\leq x\}.
\]
It is a finitary projection: it is continuous, lies below the identity,
is idempotent, and every fixed point is a directed supremum of compact
fixed points. The operation $p_{\mathrm{ty}}$ defined from $T(a)$
satisfies the three recursive equations for $\Phi$ and lies below
every continuous operation satisfying them. Consequently
$p_{\mathrm{ty}}=p$, the least fixed point defined
in \S4.1. This comparison is used for typed laws below.
Finitariness and normalization of both $p$ and all iterates $P_k$
already follow independently from
Theorem~\ref{thm:normalized-finitary}.

Filter soundness then identifies semantic typing with fixedness: for
a semantic type $a$, a value $x$ has type $a$ precisely when $p(a,x)=x$.
In particular $p(a,x)\leq x$ is in the image of this projection. We
use the usual filter interpretation $\ev M\rho$, identifying consistent
regular filters with their reflected domain elements whenever semantic
typing supplies this identification. The following compact-cover
property imposes no prescribed stage bound.

\begin{lemma}[Typed compact cover]\label{lem:typed-cover}
If a compact basis element $k$ satisfies $k\leq p(a,x)$, then there
merely exist proper compact representatives $v,c$ such that
\[
 k\leq v\leq x,\qquad c\leq a,\qquad v:c.
\]
\end{lemma}
\begin{proof}
The typed generators are closed under compatible finite joins and
compact equivalence, and include bottom. For a join, join the classifiers
as well; presupposition and conversion put both values at the common
classifier, and compatible-join typing gives the resulting generator.
Thus the generators below $x$ form a directed family with supremum
$p(a,x)$. Compactness of $k$ gives a generator above it. Replacing its
value and classifier by equivalent proper representatives preserves
the inequalities and typing. No bounded-stage interpolation is involved.
\end{proof}

\begin{lemma}[Typed relation laws]\label{lem:typed-laws}
The following implications hold for all semantic $a,x$ and syntactic
indices:
\begin{align*}
 \TyRel(p(\mathbf u,a);T,T')
  &\Longrightarrow\TyRel(p(\mathbf u,a);T',T),\\
 \TyRel(p(\mathbf u,a);T,V)\ \wedge\
 \TyRel(p(\mathbf u,a);V,W)
  &\Longrightarrow\TyRel(p(\mathbf u,a);T,W),\\
 \ElemRel(p(a,x);T,M,N)
  &\Longrightarrow\ElemRel(p(a,x);T,N,M),\\
 \ElemRel(p(a,x);T,M,N)\ \wedge\
 \ElemRel(p(a,x);T,N,O)
  &\Longrightarrow\ElemRel(p(a,x);T,M,O),\\
 \TyRel(p(\mathbf u,a);T,T')\ \wedge\
 \ElemRel(p(a,x);T,M,N)
  &\Longrightarrow\ElemRel(p(a,x);T',M,N).
\end{align*}
The relation also has head-reduction transport, expansion in both term
arguments, and left diagonal closure.
\end{lemma}
\begin{proof}
Prove the five typed implications together at finite relation depth,
with the type implications before the term implications at each
successor. The semantic projector is eliminated into lower admissible
predicates using its universe, product and function cases. In the
function case its action has the form
\[
 p(\pi(a,\varphi),f)\appl z
  =p\bigl(\varphi\appl p(a,z),\,f\appl p(a,z)\bigr).
\]
Restrict the tested argument to $p(a,z)$. Idempotence identifies the
result of testing at $z$ with that of testing at $p(a,z)$, and puts
the output in the corresponding codomain projector. The depth induction
hypotheses therefore supply the required argument conversions and output
symmetry or transitivity. The type clauses supply the head and component
evidence needed to align the tests.

Passing to the invariant requires a typed cover, not an assumption that
stage projectors preserve typing. Given a compact observation below
$p(a,x)$, Lemma~\ref{lem:typed-cover} supplies a compact value $v$
and classifier $c$ with
\[
 k\leq v\leq x,\qquad c\leq a,\qquad v:c.
\]
By Proposition~\ref{prop:projector-tower}, the compact values $v$ and $c$
are each fixed at some stage. Take the maximum of those two stages.
The value is fixed by $p(c,-)$, and the classifier is universe-typed.
Recovery transfers the premises to the finite-depth relation at this
stage; the finite-depth laws apply; recovery transfers the conclusion
back. Lowering to $k$ and compact reconstruction prove the laws on the
whole projected element. Type symmetry and transitivity use the same
argument with the universe projector.

Head transport and term expansion follow from the clauses and head
determinism. Left diagonal closure is proved first at finite depth and
then passed to the limit. It does not require semantic typing. Semantic
soundness transfers the five projected laws to typed filter denotations.
\end{proof}

\subsubsection{Fitting environments and related substitutions}

An environment fits the empty context automatically. It fits
$\Gamma,z:A$ if its restriction $\rho$ fits $\Gamma$, the denotation
$\ev A\rho$ is universe-typed, and its last value is semantically typed
by $\ev A\rho$. We write $\rho\models\Gamma$.

Let $\sigma,\sigma':\Gamma\to\Delta$ be typed substitutions. Relative
to a fitting environment $\rho$, say they are \emph{related} if for each
variable $z$ of $\Gamma$, with its type $A_z$ expressed in $\Gamma$,
we have derivations
\[
 \Delta\vdash\sigma z=\sigma'z:A_z\sigma,\qquad
 \Delta\vdash\sigma'z=\sigma z:A_z\sigma',
\]
and both relations
\begin{equation}\label{eq:related-substitutions}
 \ElemRel(\rho z;A_z\sigma,\sigma z,\sigma'z),\qquad
 \ElemRel(\rho z;A_z\sigma',\sigma z,\sigma'z).
\end{equation}
Equivalently, every proper compact observation of each raw environment
filter satisfies these relations. Typed symmetry allows reversal of
related substitutions; diagonal closure allows replacing a related pair
by its left diagonal. These statements hold in a well-formed source
context and use only the variable typings supplied by fittingness.

\begin{lemma}[Restriction to a fitting extension]\label{lem:argument-restriction}
Suppose $\rho\models\Gamma$, the substitutions $\sigma,\sigma'$ are
related, and $\Gamma\vdash A:U$. Suppose the domain induction hypothesis
provides
\[
 \TyRel(a;A\sigma,A\sigma'),\qquad a=\ev A\rho.
\]
Given an arbitrary argument test
\[
 \Delta\vdash N=N':A\sigma,\qquad
 \ElemRel(x;A\sigma,N,N'),
\]
put $y=p(a,x)$. Then $\rho[z\mapsto y]\models\Gamma,z:A$, and
$(\sigma,z\mapsto N)$ and $(\sigma',z\mapsto N')$ are related for
that extended environment.
\end{lemma}
\begin{proof}
The raw semantic projection representing $y$ is typed by $a$, since
$a$ is universe-typed. Thus the extended environment fits. Since
$y\leq x$, downward closure gives the argument relation at $A\sigma$.
The conversion law of Lemma~\ref{lem:typed-laws}, together with the type
relation at $a$, gives it at $A\sigma'$. More explicitly, the type
premise lowers to $p(\mathbf u,a)$, and idempotence identifies
$p(a,y)$ with $y$.

Syntactic cross-substitution applied to $A:U$ gives
$\Delta\vdash A\sigma=A\sigma':U$. Symmetry and conversion of the
argument equality give the required equality in the other direction at
$A\sigma'$. The old variables retain their relations, since substitution
in a weakened type ignores the newly substituted argument. These are
exactly the conditions for the two extended substitutions.
\end{proof}

\subsubsection{Why restricting the argument suffices}

For a body $M$ over $\Gamma,z:A$, let $\Lambda_{A,M,\rho}$ denote the
semantic binder filter: its finite graph entries have a compact input
$u$ typed by $\ev A\rho$ and an output observed in
$\ev M(\rho[z\mapsto u])$. This is the semantic abstraction used both
for lambda bodies and for product families.

\begin{lemma}[Evaluation through the restriction]\label{lem:binder-factor}
Let $a=\ev A\rho$ and $y=p(a,x)$. Every observation of
$\Lambda_{A,M,\rho}\appl x$ is an observation of
$\ev M(\rho[z\mapsto y])$. Hence, whenever the latter is a typed
denotation,
\begin{equation}\label{eq:binder-factor}
 \Lambda_{A,M,\rho}\appl x\leq\ev M(\rho[z\mapsto y]).
\end{equation}
The same implication holds for any function component below the binder.
\end{lemma}
\begin{proof}
For a compact input presentation $u$, consider the auxiliary family
\[
 H(u)=\ev M(\rho[z\mapsto p(a,u)]).
\]
It is monotone in the information order. At a graph entry whose input
$u$ is typed by $a$, the principal filter generated by $u$ is fixed by
$p(a,-)$. Thus that entry's output belongs to $H(u)$.

An application observation selects finitely many entries of a binder
graph under an input observed in $x$. Induction on the selection,
using monotonicity for a single entry and filter intersection for a
combined selection, puts the selected output in $H(u_0)$ at that
observed input $u_0$. Since $u_0\leq x$, projector monotonicity gives
$p(a,u_0)\leq p(a,x)$, so the output belongs to the right-hand side
of~\eqref{eq:binder-factor}. Upward closure in token order handles the
final observation. The argument is on raw filter witnesses; after
reflection, truncated witnesses are eliminated only into the desired
proposition. Compact reconstruction then transfers any lower admissible
predicate of the body denotation to the application result.
\end{proof}

There is a second point in the function and family clauses: the induction
hypothesis must supply three comparisons, all at the same left type.
Use the fitting extension from Lemma~\ref{lem:argument-restriction} and
write
\[
\begin{gathered}
 V=B(\sigma,N),\quad Y=\ev M(\rho[z\mapsto y]),\\
 P=M(\sigma,N),\quad P'=M(\sigma,N'),\quad
 Q=M'(\sigma',N),\quad Q'=M'(\sigma',N').
\end{gathered}
\]
Here $(\sigma,N)$ extends $\sigma$ by $N$ for the new variable. A self
induction hypothesis for $M$ at the diagonal base substitution yields
$\ElemRel(Y;V,P,P')$.
The cross induction hypothesis for $M,M'$ yields
\[
 \ElemRel(Y;V,P,Q'),\qquad \ElemRel(Y;V,P,Q),
\]
the second by the diagonal argument test $(N,N)$.
Typed symmetry and transitivity on $Y$ now give
\begin{equation}\label{eq:three-body-comparisons}
 \ElemRel(Y;V,P,P'),\qquad
 \ElemRel(Y;V,Q,Q'),\qquad
 \ElemRel(Y;V,P,Q').
\end{equation}
This obtains the right self comparison without changing its classifier
to $B(\sigma',N)$ and without choosing a uniform family of witnesses.

\subsection{The fundamental lemma and product inversion}

\begin{theorem}[Fundamental lemma]\label{thm:fundamental}
Suppose $\Delta$ is well formed, $\rho\models\Gamma$, and
$\sigma,\sigma':\Gamma\to\Delta$ are related relative to $\rho$.
For every derivation $\Gamma\vdash M=N:A$,
\begin{equation}\label{eq:fundamental}
\begin{aligned}
 \ElemRel(\ev M\rho;A\sigma,M\sigma,N\sigma'),\\
 \ElemRel(\ev M\rho;A\sigma',M\sigma,N\sigma').
\end{aligned}
\end{equation}
For a typing derivation $\Gamma\vdash M:A$, the same conclusions hold
with $N=M$. They hold on the whole reflected denotation and on every
compact observation below it.
\end{theorem}
\begin{proof}
Induct simultaneously on typing and equality derivations, with the
well-formed-context cases contributing no relational assertion.
The hypotheses are quantified over all fitting environments and related
substitutions. We give the cases that explain the construction.

\smallskip\noindent\emph{Variables, universe, and structural rules.}
The variable case is~\eqref{eq:related-substitutions}. The universe case
uses its own head reduction and the absence of a product reduct of $U$.
Reflexivity uses the typing hypothesis. Symmetry reverses the related
substitutions, applies the induction hypothesis, and uses typed symmetry;
semantic soundness identifies the denotations of the equal terms.
Transitivity uses the diagonal substitutions and typed transitivity.
Conversion applies the type induction hypothesis and the conversion law
in Lemma~\ref{lem:typed-laws}, at each substituted classifier.

\smallskip\noindent\emph{Dependent products.}
For formation, the premises are $\Gamma\vdash A:U$ and
$\Gamma,z:A\vdash B:U$. For equality they additionally include
$\Gamma\vdash A=A':U$ and $\Gamma,z:A\vdash B=B':U$, with the
specified typings of the bodies. Substitution and cross-substitution
give the syntactic component equalities
\[
 \Delta\vdash A\sigma=A'\sigma':U,\qquad
 \Delta,z:A\sigma\vdash B(\sigma^+)=B'(\sigma'^+):U.
\]
Here $\sigma^+$ denotes the lifted substitution under the binder.
These equalities come from the premises, not from inversion of a product
equality. The domain induction hypothesis gives the type relation on
$\ev A\rho$.

The product denotation has no universe head. Under a positive product
observation, its extracted domain is below $\ev A\rho$, and its extracted
family is below $\Lambda_{A,B,\rho}$. For an arbitrary related argument,
Lemma~\ref{lem:argument-restriction} gives a fitting extension. The body
induction hypotheses and~\eqref{eq:three-body-comparisons}, with classifier
$U$, supply the three type comparisons in~\eqref{eq:family-outputs}.
Lemma~\ref{lem:binder-factor} transfers them to the extracted family.
Fold the type clause, then the term clause at $U$. Formation is the
self-equality instance. Since the classifier is $U$, the two substituted
classifier conclusions coincide.

\smallskip\noindent\emph{Annotated abstraction.}
For $\lambda A.M:\Pi(A,B)$, its semantic value has only a function
component, so the type component of the term relation has no positive
head to test. For a related argument, restrict it as above, use the body
induction hypothesis and~\eqref{eq:three-body-comparisons}, and transfer
the outputs by~\eqref{eq:binder-factor}. Syntactic beta reduction and term
expansion replace the instantiated bodies by applications of the
substituted abstractions. This proves the three function outputs at
$\Pi(A,B)\sigma$. Apply the same argument with $\sigma,\sigma'$
reversed and then typed symmetry on the abstraction denotation to obtain
the conclusion at $\Pi(A,B)\sigma'$.

\smallskip\noindent\emph{Application and beta.}
The function relation applied to the argument relation yields the
application comparison at the left substituted result type.
Instantiating the codomain induction hypothesis under the binder gives
the relation between the two substituted result types; typed conversion
gives the right-type conclusion. This treats ordinary application and
both congruence rules. For beta, extend the environment by the argument
denotation and the substitutions by the two substituted arguments.
The body hypothesis applies there. Semantic beta and substitution
identify the denotation, while syntactic beta and term expansion give
the required equality comparison at both result types.

\smallskip\noindent\emph{Typed function extensionality.}
Suppose the premises include
\[
 \Gamma,z:A\vdash f^{\uparrow}z=g^{\uparrow}z:B,
 \qquad \Gamma\vdash f:\Pi(A,B),\quad
 \Gamma\vdash g:\Pi(A,B).
\]
Semantic typing gives the eta inclusions
\begin{equation}\label{eq:semantic-eta}
 \ev f\rho\leq\Lambda_{A,f^{\uparrow}z,\rho},\qquad
 \ev g\rho\leq\Lambda_{A,g^{\uparrow}z,\rho}.
\end{equation}
To see the first, take a typed compact cover of an observation of
$\ev f\rho$. A non-bottom cover must be a function graph: universe and
product values cannot have a product classifier. Each graph entry has a
key typed by $\ev A\rho$ and its output is observed in the application
of $\ev f\rho$ to that key. Semantic weakening identifies this with the
corresponding eta body. All entries therefore belong to the binder,
and the graph covers the original observation. The proof for $g$ is
the same. This is a consequence of semantic typing and soundness.

The pointwise equality induction hypothesis, together with its symmetry
and transitivity consequences, supplies the self and cross hypotheses
for the two eta bodies. Restrict each tested argument to a fitting
extension, obtain~\eqref{eq:three-body-comparisons}, and use
\eqref{eq:semantic-eta} and Lemma~\ref{lem:binder-factor} to transfer them
to applications of $\ev f\rho$. Substitution in $f^{\uparrow}z$ is exactly
application of $f\sigma$ to the substituted argument. Hence the three
outputs fold to the required function relation. For the right classifier,
reverse both the pointwise equality and the substitutions, use typed
symmetry, and transfer from the denotation of $g$ to that of $f$ by
semantic soundness of the extensionality rule.

These cases cover every derivation constructor. The proof can be read
observation by observation on the raw filters. Semantic typing gives
consistent regular denotations, and Lemma~\ref{lem:reconstruction}
then supplies the whole-denotation assertions. Conversely, downward
closure gives every compact-observation conclusion.
\end{proof}

\begin{corollary}[Mere product inversion]\label{cor:independent-pi}
From a derivation $\Gamma\vdash\Pi(A,B)=\Pi(A',B'):T$, for any
classifier $T$, one obtains
\[
 \trunc{(\Gamma\vdash A=A':U)\mathbin{\times}
          (\Gamma,z:A\vdash B=B':U)}.
\]
\end{corollary}
\begin{proof}
The environment $\rho_\bot$ assigning bottom to every variable fits
every well-formed context. Prove this by induction on the context:
soundness makes the next type universe-typed, and bottom is a value of
that type. The identity substitution is related to itself at $\rho_\bot$
by strictness of the relation.

The literal product denotation has the compact observation
$k_\Pi=\pi(\bot,\bot)$: use the bottom observation of its domain and the
empty family graph. Apply Theorem~\ref{thm:fundamental} at the identity
substitutions and observe $k_\Pi$. The type component of the term clause
gives $\TyRel(k_\Pi;\Pi(A,B),\Pi(A',B'))$. Unfolding the positive
product clause supplies a mere tuple of evidenced heads with their
component equalities. Literal products are head normal, so head
determinism identifies those components with $A,B,A',B'$. Mapping this
truncated tuple to its equalities gives the displayed conclusion.
\end{proof}

The argument derives product inversion after the fundamental lemma; it
uses neither syntactic product injectivity nor an earlier adequacy
theorem as a premise. The conclusion is explicitly truncated. It does
not assert a chosen pair of derivation trees, or a simultaneous choice
of witnesses over all compact observations.

\clearpage
\appendix
\section{Embedding--projection categories and initial algebras}\label{app:categories}

In this appendix we also use an \emph{unpresented} Scott domain: its small
compact basis is merely known to exist, rather than selected as object
data. A presentation supplies such an object by propositional truncation.
\begin{proposition}[Univalence of Scott-domain objects]
Let $D$ and $E$ be unpresented Scott domains. The canonical map is an
equivalence
\[
 (D=E)\simeq\{\text{order isomorphisms between $D$ and $E$}\}.
\]
\end{proposition}
\begin{proof}
Suprema are unique, so completeness contributes no choices for a fixed
order. Least elements are unique, and the basis-existence condition is a
proposition. It therefore suffices to compare orders. Fix an order $P$.
The total space of orders equipped with an order isomorphism from $P$ is
contractible: univalence makes the type of carrier equivalences from
$P$ contractible, while propositional univalence and the order laws make
the relation over each such equivalence unique. The total space for Scott
domains is a retract of this space equipped with the remaining structure
in proposition-valued fibres; it is contractible as well. Projecting its
unique path from $(D,\mathrm{id})$ to $(E,i)$ produces the inverse of the
canonical map. Path induction and contractibility give both inverse laws.
\end{proof}

\begin{proposition}[Univalent categories of Scott domains]
Unpresented Scott domains form a univalent category with Scott-continuous
maps. They also form a univalent category whose morphisms $D\to E$ are
embedding--projection pairs of continuous maps
$e:D\to E$ and $p:E\to D$ satisfying $p\circ e=\mathrm{id}_D$ and
$e\circ p\leq\mathrm{id}_E$ pointwise. In both categories, categorical
isomorphisms correspond to order isomorphisms of domains.
\end{proposition}
\begin{proof}
Identity and composition preserve continuity. For pairs, composition is
$(e_2\circ e_1,p_1\circ p_2)$; monotonicity and the two pair laws give its
retraction and pointwise inequality. An invertible continuous map has a
continuous inverse: an order isomorphism preserves and reflects directed
suprema. Conversely, the underlying map of a categorical isomorphism and
its inverse give an order isomorphism. For an invertible pair, its
projection agrees with the inverse of its embedding, so its entire arrow
is determined by that order isomorphism. Thus the singleton total space
of categorical isomorphisms from a fixed domain is a retract of the
proposition-valued total space used in the preceding result. The
total-space criterion for univalence applies to both precategories.
\end{proof}

\begin{theorem}[Compact elements and ideal completion]
\label{thm:cusl-equivalence}
Let $\mathbf{Dom}_{\triangleleft}$ be the univalent category of unpresented
Scott domains and continuous embedding--projection pairs. Let
$\mathbf{CUSL}$ be the univalent category of conditional upper
semilattices with merely small carrier. Its morphisms preserve bottom and
compatible binary joins and reflect both order and compatibility. Then
\[
  \mathbf{Dom}_{\triangleleft}\simeq\mathbf{CUSL}.
\]
The functors are $\K(D)$, the CUSL of compact elements, and
$\operatorname{Idl}(P)$, the Scott domain of ideals of $P$. There are natural
isomorphisms
\[
  P\cong\K(\operatorname{Idl}(P)),\qquad
  D\cong\operatorname{Idl}(\K(D)).
\]
\end{theorem}
\begin{proof}[Proof outline]
For a small presentation of $P$, ideals are inhabited, downward closed,
directed lower subsets. Inclusion is a small relation. Directed suprema
are unions; the supremum of a bounded family is the ideal generated by
its finite joins. Principal ideals are compact and form a small basis.
An embedding of CUSLs extends to the downward closure of its direct image;
inverse image supplies the continuous projection. The four embedding
laws give the retraction and deflation equations.

Conversely, an embedding--projection pair sends compact elements to
compact elements: its embedding is left adjoint to its projection, and the
projection preserves directed suprema. The restricted embedding preserves
compatible joins and reflects compatibility. A small basis presents the
compact-element carrier via the replacement image; the presentation is
only asserted to exist, so it does not become object data.
For a set, univalence makes two small presentations equal. Thus mere
smallness supplies the small presentation used for ideal completion
without a choice axiom. Every compact ideal is principal, while each
domain element is the supremum of the compact elements below it. These
give the displayed isomorphisms. Their naturality is checked on principal
ideals and on the projections of embedding--projection pairs, respectively.
\end{proof}

\subsection{Countable chains of presented domains}

\begin{proposition}[Countable limits of presented embedding--projection chains]
Let $D_0\to D_1\to\cdots$ be a chain of presented Scott domains and
continuous embedding--projection pairs $(e_n,p_n)$. The coherent sequences
\[
 L=\{(x_n)_n\mid p_n(x_{n+1})=x_n\}
\]
form a presented Scott domain. Each $D_n$ embeds into $L$, with projection
given by the $n$th coordinate, and every $x\in L$ is the directed supremum
of its stage approximants $e_{n,\infty}(x_n)$. The compact basis consists of
stage lifts of compact basis elements. This cone is a colimit in the
embedding--projection category on presented objects: every compatible
cone to a presented Scott domain has a unique mediating pair.
\end{proposition}
\begin{proof}
Orders and directed suprema of coherent sequences are coordinatewise.
For bounded suprema, take the supremum at each stage, lift it to $L$,
and then take the directed supremum of these stage lifts. A stage lift is
compact when its stage element is compact. Stage approximation and compact
density at each stage give density in $L$ without countable choice.

For a compatible cone $(u_n,v_n):D_n\to Y$, the mediating projection sends
$y$ to $(v_n(y))_n$; its embedding sends $(x_n)_n$ to
$\lub_n u_n(x_n)$. The cone equations make the latter family increasing.
The retraction and deflation laws follow by checking each stage coordinate
and using continuity. A second mediating pair agrees on each stage lift,
so stage exhaustion and continuity make its embedding equal to this one;
its projection is equal coordinatewise.

\end{proof}

\subsection{Initial algebras and categorical transfer}

\begin{theorem}[Abstract initial-chain theorem]\label{thm:initial-chain}
Let $\mathcal C$ be a category with an initial object and colimits of
countable chains. If an endofunctor $F:\mathcal C\to\mathcal C$ preserves
these colimits, then the colimit $\mu F$ of
\[
 0\longrightarrow F0\longrightarrow F^2 0\longrightarrow\cdots
\]
has an initial $F$-algebra structure
$\mathsf{fold}:F(\mu F)\to\mu F$. The map $\mathsf{fold}$ is an
isomorphism. If $\mathcal C$ is univalent, the isomorphism induces an
equality $\mu F=F(\mu F)$.
\end{theorem}
\begin{proof}
Removing the initial object from the chain leaves a chain with the same
colimit: a cocone over the tail extends uniquely at stage zero. Preservation
by $F$ makes $F(\mu F)$ another colimit of this tail. The unique mediating
arrows between these two colimits give fold and unfold; their composites
are identities by colimit uniqueness.

For an algebra $a:F(Y)\to Y$, define its stage maps recursively by the
unique arrow $0\to Y$ and $h_{n+1}=a\circ F(h_n)$. The functor laws make
them a cocone. Its mediator $h:\mu F\to Y$ satisfies
$h\circ\mathsf{fold}=a\circ F(h)$ because both sides agree on every
image stage and $F(\mu F)$ is their colimit. Any other algebra morphism
agrees with the stage maps by induction, hence equals $h$ by colimit
uniqueness. Since hom types are sets, the algebra equation is a
proposition, so the type of algebra morphisms is contractible.
\end{proof}

The theorem is formalized for precategories, with a separate generic
univalent-category corollary for the object equality.

\begin{theorem}[Countable CUSL colimits and initial algebras]
\label{thm:cusl-initial}
The univalent category $\mathbf{CUSL}$ of conditional upper semilattices
and embeddings has an initial object and colimits of countable chains.
For any endofunctor $F:\mathbf{CUSL}\to\mathbf{CUSL}$ preserving these
colimits, there is an initial $F$-algebra
$\mathsf{fold}:F(\mu F)\to\mu F$. Its structure map is an isomorphism,
and univalence gives $\mu F=F(\mu F)$. The preservation hypothesis is a
proposition.
\end{theorem}
\begin{proof}
The one-point CUSL is initial. For a chain of CUSL embeddings, form the
chain of their presented ideal domains and its coherent-sequence EP limit.
Take the compact elements of this limit as the colimit CUSL. The unit of
the CUSL--EP equivalence and the canonical EP cone give the cocone legs.
For a cocone into $Y$, ideal extension gives an EP cocone into the presented
ideal domain of $Y$; its universal mediator restricts to a CUSL embedding
from the compact elements of the limit to $Y$. Naturality of the unit gives
the factorization equations. Every compact element of the limit is equal
to the lift of a compact basis element at some stage, so agreement on the
cocone legs determines the mediator uniquely. The abstract initial-chain
theorem now supplies the initial algebra and the equality.

For a fixed cocone, two mediators satisfying the universal property agree
by uniqueness. All remaining laws are equalities in hom sets, so its
colimit structure is a proposition. Function extensionality then shows
that preservation of all existing countable-chain colimits is a
proposition. The small presentations used in ideal completion are
selected by univalence from mere smallness; no countable choice is used.
\end{proof}

\begin{theorem}[Categorical structure identity principle]
\label{thm:category-sip}
At the object and hom universe levels used here, an equivalence between
univalent categories $\mathcal C$ and $\mathcal D$ induces an equality
$\mathcal C=\mathcal D$ of category records. In particular,
$\mathbf{CUSL}=\mathbf{Dom}_{\triangleleft}$.
\end{theorem}
\begin{proof}
This is the equivalence-to-identity direction of the categorical
structure identity principle of
\href{https://arxiv.org/abs/1303.0584}{Ahrens--Kapulkin--Shulman}.
The two natural reconstruction isomorphisms give mutually inverse
functions on object types after category univalence converts their
components to object equalities. Univalence of the object universe then
identifies the object types. The functor induces an equivalence on each
hom type: its inverse is obtained by applying the quasi-inverse functor
and conjugating by the reconstruction isomorphisms. Fiberwise univalence
identifies the hom families. The functor laws identify identities and
composition after transport. Because hom types are sets, the remaining
category laws are propositions; univalence of a fixed precategory is
also a proposition. The resulting equality of precategories extends to
equality of category records. The checked construction is
\texttt{CategorySIPTheorem.On.equality}; it uses only the previously
stated universe univalence and function extensionality.
\end{proof}

\begin{corollary}[Unpresented Scott EP initial algebras]
\label{cor:ep-initial}
The category of unpresented Scott domains and embedding--projection pairs
has colimits of countable chains. Every endofunctor of this category that
preserves these colimits has an initial algebra with an invertible
structure map and an object equality $X=F(X)$.
\end{corollary}
\begin{proof}[Proof using univalence]
By Theorem~\ref{thm:category-sip}, the CUSL--EP equivalence gives an
equality of categories. Transport the CUSL countable-chain colimit
structure along this equality; this is the definition
\texttt{ScottEPTransport.cocomplete}. The endofunctor $F$ remains an
endofunctor of EP throughout: only the colimit structure is transported.
The one-point EP domain is initial, so
Theorem~\ref{thm:initial-chain} applies. Univalence turns its solution
isomorphism into equality of objects.
\end{proof}

\begin{proof}[Proof up to isomorphism using a precategory equivalence]
The categorical argument needs only an equivalence of precategories
$I:\mathbf{CUSL}\rightleftarrows\mathbf{Dom}_{\triangleleft}:K$
with natural reconstruction isomorphisms, together with the CUSL chain
colimits. It does not need an equality of the two category records.
For an EP chain $S$, take a colimit $(L,\lambda_n)$ of the chain $KS$.
The EP colimit vertex is $I L$, with legs
$I(\lambda_n)\circ\eta_{S_n}:S_n\to I L$, where
$\eta_{S_n}:S_n\cong IK(S_n)$ is the reconstruction isomorphism.
Given a cocone $q_n:S_n\to Y$, the arrows $K(q_n)$ have a unique
mediator $m:L\to K Y$. Compose $I(m)$ with the inverse of
$\eta_Y:Y\cong IKY$ to obtain a mediator $IL\to Y$.
Naturality of $\eta$ proves the factorization equations; full
faithfulness of $I$ proves uniqueness. Thus EP has countable-chain
colimits as a precategory. Apply Theorem~\ref{thm:initial-chain} to the
initial EP object and an endofunctor $F$ preserving these colimits.
Its initial-chain colimit $X$ has inverse EP arrows
$\mathsf{fold}:F(X)\to X$ and $\mathsf{unfold}:X\to F(X)$, so
$X\cong F(X)$. Here continuity of $F$ means preservation of
countable-chain colimits. Function and propositional extensionality suffice for
the equalities of maps and proof fields in this argument. If objects
are passed to an isomorphism-class quotient, it also gives
$[X]=[F(X)]$ there; an equality $X=F(X)$ of domain objects uses the
univalence step in the preceding proof.
\end{proof}

\begin{remark}[Scope of the formalization without univalence]
The precategory transfer is checked abstractly in
\texttt{CategoryPreEquivalenceColimits}; the initial-chain
construction is checked for arbitrary precategories in
\texttt{CategoryInitialChain}.

The current concrete construction of $I$ and of the CUSL colimits is
inside the development parameterized by univalence. In particular,
\texttt{ScottCUSLIdeals.Canonical} calls
\texttt{ScottSmallness.choose} to extract a small presentation from
mere smallness; that operation uses univalence of the small universe.
The checked concrete transfer to presented EP also uses this canonical
completion. Consequently, the abstract second proof is formalized,
but the full concrete corollary for unpresented Scott EP has not been
separately checked from only propositional extensionality, function
extensionality, and quotients. A construction
of the needed ideal completion from those assumptions alone remains
to be supplied.
\end{remark}

\begin{corollary}[Transport of a CUSL domain equation]
\label{cor:cusl-solution-transport}
Let $F:\mathbf{CUSL}\to\mathbf{CUSL}$ preserve countable-chain
colimits. Under the equality
$\mathbf{CUSL}=\mathbf{Dom}_{\triangleleft}$ of
Theorem~\ref{thm:category-sip}, transport of $F$ and its initial algebra
gives an EP endofunctor $F^{\mathrm{EP}}$ with an initial algebra
$\mathsf{fold}^{\mathrm{EP}}:F^{\mathrm{EP}}(\mu^{\mathrm{EP}}F)
\to\mu^{\mathrm{EP}}F$. Its fold map is invertible and
$\mu^{\mathrm{EP}}F=F^{\mathrm{EP}}(\mu^{\mathrm{EP}}F)$.
\end{corollary}
\begin{proof}
Theorem~\ref{thm:cusl-initial} gives a dependent package comprising
$F$, its colimit-preservation proof, the initial $F$-algebra, invertibility
of its structure map, and the object equality. Transport this whole
package along the category equality. The resulting package has exactly
the stated fields for an EP endofunctor. The checked construction is
\texttt{ScottCUSLSolutionTransport.On}; its algebra-morphism center and
uniqueness fields establish initiality after transport.
\end{proof}

\begin{theorem}[Univalence of the algebra category]
For any univalent category $\mathcal C$ and endofunctor
$F:\mathcal C\to\mathcal C$, the category $\operatorname{Alg}(F)$ is
univalent. Its objects are pairs $(X,a:F(X)\to X)$, and its arrows
$h:(X,a)\to(Y,b)$ satisfy $h\circ a=b\circ F(h)$.
\end{theorem}
\begin{proof}
An algebra isomorphism gives an isomorphism of underlying objects, hence
an equality $X=Y$ by univalence of $\mathcal C$. The commuting square
identifies the transported algebra map with $b$, giving equality of the
algebras. Path induction and the sethood of hom types show that this
construction and the canonical map from equalities to isomorphisms are
inverse. The functor laws give identities and composition of algebra
arrows. No colimit preservation assumption on $F$ is used.
\end{proof}

The direct coherent-sequence construction uses \emph{presented} Scott
domains. Transport by category equality extends countable-chain colimits
to arbitrary unpresented objects with merely existing compact bases.
Here ``$\omega$-continuous'' means
preservation of existing countable-chain colimits. An implication from
continuity on ordered hom types to this preservation is not asserted.
The preservation property is a proposition. A separate generic proof
transports countable-chain colimits directly along a categorical
equivalence, without using equality of categories; the corollary above
uses the stronger structure identity theorem.

\clearpage
\section*{Formal checks and total timing}\phantomsection\label{sec:formal}

The development is checked in Unit with univalence and replacement as
ambient parameters. Function extensionality is derived from univalence;
propositional truncation is obtained from the replacement interface.
No excluded middle, choice or propositional resizing axiom is added.
The results discussed here are drawn from the following checked developments:
\begin{enumerate}[leftmargin=*,itemsep=3pt]
\item Scott domains with small compact bases, including a univalent object
type with merely existing basis and univalent categories of continuous
maps and embedding--projection pairs; continuous function spaces,
compact step approximations, least fixed points, and the abstract
initial-chain theorem, Theorem~\ref{thm:initial-chain}, and the generic
univalent category of endofunctor algebras; the equivalence of univalent
Scott EP and CUSL categories, Theorem~\ref{thm:cusl-equivalence}.
The CUSL initial-algebra theorem, Theorem~\ref{thm:cusl-initial},
uses the checked countable-chain colimit construction; the categorical
structure identity principle, Theorem~\ref{thm:category-sip}, transfers it to
the unpresented EP corollary, Corollary~\ref{cor:ep-initial}, and directly
transports each CUSL solution by Corollary~\ref{cor:cusl-solution-transport}.
\item The Scott filter domain and recursive projector, with independent
preservation of normalized finitary families and their Kleene limit
(Theorem~\ref{thm:normalized-finitary}); the finite-stage projectors and their structural laws, finite stabilisation, and comparison
with bounded compact typing. The existing Unit proof constructs these
projectors using ranks of finite presentations. It also checks the
recursive successor equation used in
Proposition~\ref{prop:projector-tower}. A proof of all the tower laws
starting solely from that recursive definition has not been separately
formalised.
\item The independent strict lower admissible invariant, its equations
and uniqueness; typed symmetry, transitivity and conversion.
\item Typed argument restriction; Pi formation and equality, annotated
lambda and function extensionality; the complete derivation induction.
\item Whole-denotation adequacy at both substituted types, compact
recovery, and mere syntactic Pi inversion.
\end{enumerate}
The complete induction is \texttt{At.fundamental} in
\path{IndependentFundamental.seg}; its consequences are in
\path{IndependentAdequacy.seg} and \path{IndependentInversion.seg}.

On 13 November 2026 the updated \texttt{python3 UNIT/check.py} passed
the full development (14,856 declarations), the general dependent-product
and dependent-sum session (1,156), the independent projector session (3,838),
and the independent logical-relation session (6,338), each with zero failures.
The general dependent-product session has no filter-domain or typing
modules in scope. The projector session excludes the semantic projector,
typed stage comparisons, and the older direct common-code proof. Both import boundaries are audited by the script.
All 40 dependency hashes remain unchanged. The previously recorded timing
below concerns the older two-session check.

On 14 November 2026, after adding admissible fixed-point induction and
using it for the projector invariant, the same script passed all four
sessions: 14,883, 1,156, 3,849, and 6,349 declarations respectively,
each with zero failures. All 40 frozen dependency hashes matched.

After the subsequent general CUSL construction of $\mathrm{Fp}$ on
14 November, the script passed five sessions: the full development
(15,724 declarations), the independent CUSL $\mathrm{Fp}$ session
(2,412), the general product/sum session (1,156), the independent
projector session (3,849), and the independent logical-relation session
(6,349), each with zero failures. The CUSL session checks the general
functor, small directed-colimit preservation, adjunction and action
formulas with the filter and typing modules out of scope. Its import
boundary is audited by the script. All 40 dependency hashes matched.

All five sessions also pass with a checker built from the published Unit
revision \texttt{18bb219f}. Reproduction instructions and the full revision
identifier are in \path{UNIT/README.txt}.

A run on 12 November 2026 of
\texttt{/usr/bin/time -p python3 UNIT/check.py}, using the recorded local
macOS arm64 checker, gave the following results. The type-checker is a
compiled Haskell program. Python only verifies the dependency hashes and
invokes that program for the two checks; \texttt{time -p} measures the
complete command.
\begin{center}
\begin{tabular}{lrr}
\hline
Check & Declarations & Checker CPU time\\
\hline
Complete development & 12,187 & 8.56 s\\
Independent relation and consequences & 4,909 & 1.18 s\\
\hline
\end{tabular}
\end{center}
The \textbf{total wall-clock time was 10.46 seconds}, including dependency
verification and both checker processes. The independent run has a
104-module dependency closure with the old logical relation and the earlier
proofs of adequacy, syntactic Pi injectivity and subject reduction out of scope. All 40 frozen
dependency hashes were verified unchanged. The two declaration counts
overlap; the full count also includes developments outside this account.
These timings describe this local run, not a machine-independent bound.

For comparison, on 14 November 2026 we also checked the original MIN
proof of domain-model soundness, logical-relation adequacy,
$\Pi$-injectivity and subject reduction. Full-source checks took \textbf{0.54 seconds
wall-clock time in Unit} (1,800 declarations, zero failures; 0.38 seconds
checker CPU time) and \textbf{26.42 seconds in Agda 2.8.0} (54 modules).
The Agda command used \texttt{--ignore-interfaces --without-K} on
\path{MIN/SubjectReduction.agda}; Unit checked the original
\path{domain/mltt/All.seg}, including both versions of its logical
relation. These sequential runs concern the two implementations of the
earlier finite-element proof; the larger Scott-domain development timed
above also contains the additional constructions described in this note.
Commands and measurements are recorded in \path{UNIT/MIN_TIMINGS.txt}.

The inversion conclusion remains propositionally truncated. Comparison
with the old uniform proof-relevant relation and transfer to the stable
reflexive domain remain separate. The checked projector tower has not
been identified with the separately constructed CUSL initial-chain
projections; that comparison requires finite-stage coverage and
compatibility with application. The classical finite-poset
counterexample is illustrative and has not been formally checked.

\begin{thebibliography}{9}
\bibitem{cgw}
T. Coquand, C. Gunter and G. Winskel,
\emph{Domain theoretic models of polymorphism},
Information and Computation 81 (1989), 123--167.

\bibitem{coquand-huber}
T. Coquand and S. Huber,
\emph{An adequacy theorem for dependent type theory},
Theory of Computing Systems 63 (2019), 647--665.
\href{https://doi.org/10.1007/s00224-018-9879-9}{doi:10.1007/s00224-018-9879-9}.

\bibitem{dejong-escardo}
T. de Jong and M. H. Escard\'o,
\emph{Continuous and algebraic domains in univalent foundations},
Journal of Pure and Applied Algebra 229 (2025), 108072.
\href{https://doi.org/10.1016/j.jpaa.2025.108072}{doi:10.1016/j.jpaa.2025.108072}.
\end{thebibliography}

\end{document}
